\renewcommand{\Anno}{2003}

\renewcommand{\NumeroProblema}{2}

\renewcommand{\Testo}{Sia $P$ un punto interno ad una circonferenza
$\Gamma$, distinto dal centro $O$.  Determinare i punti $Q$ di
$\Gamma$ per cui l'angolo $O\widehat{Q}P$ risulta massimo.}

\renewcommand{\Risposta}{Sono i due punti $Q$ della circonferenza per
cui $QP$ è perpendicolare ad $OP$.}

\renewcommand{\Aiutino}{Vi sono almeno due approcci possibili.
\begin{itemize}
	\item Per via trigonometrica, applicando il teorema dei seni al triangolo $OPQ$, in modo da ricavare (il seno del)l'angolo incognito in funzione di $OP$ e $OQ$ (noti) e del seno dell'angolo $O\widehat{P}Q$.

	\item  Per via sintetica, mostrando preliminarmente che nelle configurazioni ottimali la circonferenza circoscritta ad $OPQ$ deve essere tangente a $\Gamma$.
\end{itemize}}

\renewcommand{\SoluzioneUfficiale}{
}

\renewcommand{\SoluzioniAlternative}{
}

\renewcommand{\Osservazioni}{
}

\renewcommand{\Commenti}{
}
% Aggiunta redazionale: la soluzione preesistente non viene modificata.
\providecommand{\SoluzioneRedazionale}{}
\renewcommand{\SoluzioneRedazionale}{
% SSSUP-BOZZA-AI: SoluzioneRedazionale; autore=ChatGPT; data=2026-11-08; stato=DA-VERIFICARE
% SSSUP-PROVENIENZA: AI: elaborazione nuova dal testo
Siano $r=OQ$ il raggio della circonferenza e $d=OP$, con $0<d<r$. Poniamo $\alpha=\angle OQP$ e $\beta=\angle OPQ$. Dal teorema dei seni nel triangolo $OPQ$ segue
\[
\sin\alpha=\frac d r\sin\beta\leq\frac d r.
\]
L'angolo $\alpha$ è sempre acuto: se fosse retto o ottuso, il suo lato opposto $OP$ dovrebbe essere più lungo di $OQ$, mentre $d<r$. Massimizzare $\alpha$ equivale quindi a massimizzarne il seno. L'uguaglianza si ha esattamente per $\beta=\pi/2$, ossia quando $QP$ è perpendicolare a $OP$.

La perpendicolare a $OP$ per $P$ incontra la circonferenza in due punti, che sono entrambi e soltanto i massimizzatori. Il massimo angolo vale $\arcsin(d/r)$.
}
